% !TEX program = lualatex
% !TeX encoding = UTF-8
\PassOptionsToPackage{unicode}{hyperref}
\PassOptionsToPackage{bookmarksnumbered,bookmarksopen}{hyperref}
\documentclass[12pt,a4paper]{article}

% ============================================================
% 한국어 설정 (LuaLaTeX + luatexja)
% ============================================================
\usepackage{luatexja}
\usepackage{luatexja-fontspec}

% 한국어 고딕체 (sans) – Noto Sans KR
\setsansjfont{Noto Sans KR}[
UprightFont = * Regular,
BoldFont = * Bold,
Script = Hangul,
Language = Korean
]

% 한국어 명조체 (serif) – Noto Serif KR
\IfFontExistsTF{Noto Serif KR}{
	\setmainjfont{Noto Serif KR}[
	UprightFont = * Regular,
	BoldFont = * Bold,
	Script = Hangul,
	Language = Korean
	]
}{
	% fallback
	\setmainjfont{Noto Sans KR}[
	UprightFont = * Regular,
	BoldFont = * Bold,
	Script = Hangul,
	Language = Korean
	]
}

% 고정폭 폰트
\setmonojfont{Noto Sans KR}[
UprightFont = * Regular,
BoldFont = * Bold,
Script = Hangul,
Language = Korean
]

% 라틴 문자 폰트
\setmainfont{Times New Roman}[Ligatures=TeX]
\setsansfont{Arial}[Ligatures=TeX]
\setmonofont{Latin Modern Mono}[
WordSpace={1,0,0},
HyphenChar=None,
PunctuationSpace=WordSpace
]

% ============================================================
% 기타 패키지
% ============================================================
\usepackage{amsmath,amssymb,amsthm}
\usepackage{graphicx}
\usepackage{hyperref}
\usepackage{url}
\usepackage{xcolor}
\usepackage{booktabs}
\usepackage{geometry}
\geometry{margin=2.5cm}

% ============================================================
% YUCT 매크로
% ============================================================
\newcommand{\Keff}{K_{\mathrm{eff}}}
\newcommand{\kappac}{\kappa_c}
\newcommand{\eps}{\varepsilon}
\newcommand{\betagamma}{\beta\gamma}
\newcommand{\betac}{\beta}
\newcommand{\tauf}{\tau}

% ============================================================
% 정리 환경 (한국어)
% ============================================================
\newtheorem{theorem}{정리}
\newtheorem{lemma}{보조정리}
\newtheorem{corollary}{따름정리}
\newtheorem{definition}{정의}
\newtheorem{prediction}{예측}

% ============================================================
% 하이퍼링크 설정
% ============================================================
\hypersetup{
	colorlinks=true,
	linkcolor=blue,
	citecolor=blue,
	urlcolor=blue,
	pdftitle={부록 S: 냉각 원자에서 광자의 음의 시간 지연 — 조정된 양자 역학의 발현: Yakushev 통합 조정 이론 (YUCT) 내에서의 수정된 분석},
	pdfauthor={Alexey V. Yakushev},
	pdfkeywords={YUCT, 조정 효율, 음의 군 지연, 광자-원자 상호작용, 로그 스케일링, 온도 독립성, 철회}
}

\begin{document}
	
	\begin{titlepage}
		\centering
		\vspace*{1cm}
		{\Huge\bfseries 부록 S: 냉각 원자에서 광자의 음의 시간 지연 — 조정된 양자 역학의 발현: \\ Yakushev 통합 조정 이론 (YUCT) 내에서의 수정된 분석}
		\vspace{1.5cm}
		
		{\large Alexey V. Yakushev\textsuperscript{1}}
		\vspace{0.3cm}
		{\normalsize \textsuperscript{1}Yakushev Research, \\ 
			\textbf YUCT 이론: \url{https://yuct.org/}, \\ 
			\textbf YPSDC 프로토콜: \url{https://ypsdc.com/}\\
			\textbf DOI: \url{https://doi.org/10.5281/zenodo.18444598}}
		\vspace{1cm}
		{\normalsize 2026년 5월 \\ \large 버전 V.3.0 \\(수정된 로그 스케일링 및 13\% 온도 예측의 철회)}
		
		\vfill
		
		\begin{abstract}
			\noindent
			최근의 획기적인 실험 (Angulo 등, 2024)은 극저온 원자 구름을 통과하는 광자가 음의 시간 지연으로 매질을 떠날 수 있음을 보여주었다 – 파동 묶음의 피크가 들어간 것보다 더 일찍 나타난다. 이 반직관적인 현상은 인과율을 위반하지 않으면서도 깊은 양자 간섭 효과를 드러낸다. 본 연구에서 우리는 Yakushev 통합 조정 이론 (YUCT) 내에서의 설명을 재검토한다. 범용 스케일링 법칙 \(\eps = \kappac \alpha (\ln\Keff)^{\betac}\) (\(\betac = 2/3\), \(\kappac = 1/3\), 부록 X에서 유도)에 기초하여, 군 지연 \(\tau_g\)와 조정 효율 \(\Keff\) 사이의 올바른 관계가 로그적임을 보여준다: \(|\tau_g| = \alpha_0 \ln\Keff\). 음의 시간 지연의 존재는 이미 실험에 의해 확인되었으며; \((-1.14 \pm 0.22)\tau_0\)까지의 값이 관찰되었다. 그러나 이전의 선형 가정과 달리, 로그 형태는 \(\tau_g\)의 온도 의존성이 극도로 약함을 의미한다. 표준 온도 저하 관계 \(\Keff(T) = K_0 (T/T_0)^{-\gamma}\) (\(\gamma \approx 0.30\), 지구물리학적 데이터에서 유도)를 사용하여, 온도를 두 배로 할 때 \(|\tau_g|\)의 상대적 변화가 \(10^{-8}\) 정도임을 발견하며, 이는 현재 실험의 감도보다 훨씬 낮다. 따라서 이전에 주장된 13\% 효과는 잘못된 선형 모델의 인공물이며 철회되어야 한다. 우리는 수정된 이론적 틀을 제공하고, 광학 깊이나 펄스 지속 시간의 변화와 같이 로그 스케일링에 민감할 수 있는 대체 실험적 검증을 논의한다. YUCT의 핵심 예측 – 음의 군 지연이 조정 효율의 발현이라는 것 –은 여전히 유효하지만, 냉각 원자 시스템에서의 검증 가능성은 개선된 설정을 요구한다.
		\end{abstract}
		
		\vspace{0.5cm}
		\noindent\textbf{키워드:} YUCT, 조정 효율, 음의 군 지연, 광자-원자 상호작용, 로그 스케일링, 온도 독립성, 철회.
		
	\end{titlepage}
	
	\tableofcontents
	\newpage
	
	\section{서론}
	
	공명 매질을 통한 빛의 전파는 수십 년 동안 양자 광학의 초석이었다. 군 지연, 펄스 재형성, 초광속 효과는 이론적 및 실험적으로 잘 이해되어 있다 \cite{Brillouin1960, Milonni2004}. 그럼에도 불구하고, 근본적인 질문이 남아 있다: 광자가 원자 여기로 보내는 시간의 물리적 의미는 무엇인가?
	
	Angulo 등 \cite{Angulo2024}의 획기적인 실험은 투과된 광자가 별도의 탐침 빔에 유도하는 교차-커 위상 변이를 측정함으로써 이 문제를 다루었다. 그들은 투과된 광자에 대한 적분 원자 여기 시간 \(\tau_T\)가 음수가 될 수 있으며, \(\tau_g < 0\)일 때조차 군 지연 \(\tau_g\)와 일치함을 발견했다. 구체적으로, 실험은 \(\tau_T/\tau_0 = -1.14 \pm 0.22\) (10 ns 펄스, OD~4) 및 \(0.54 \pm 0.28\) (다른 매개변수)의 값을 보고했다 \cite{Angulo2024, Nixon2024}. 이 놀라운 결과는 단순한 직관에 도전하며 더 깊은 이론적 틀을 요구한다.
	
	Yakushev 통합 조정 이론 (YUCT) \cite{Yakushev2026a}은 그러한 틀을 제공한다. YUCT는 양자 광학에서 중력에 이르기까지 모든 물리적 상호작용이 \textbf{조정 효율} \(\Keff\)에 의해 정량화되는 근본적인 조정 과정의 발현이라고 가정한다. 범용 스케일링 법칙은 임의의 상대 오차 \(\eps\)를 \(\Keff\)와 연결한다 (부록 X, 식~\ref{eq:error-law-corrected} 참조):
	\begin{equation}
		\eps = \kappac \alpha (\ln\Keff)^{\betac}, \qquad \betac = \frac{2}{3},\ \kappac = \frac{1}{3},
		\label{eq:scaling}
	\end{equation}
	여기서 \(\alpha\)는 시스템 의존적이다. 이 법칙은 DNA 복제 오차에서 우주 마이크로파 배경 변동까지 40 자릿수에 걸쳐 검증되었다 \cite{Yakushev2026c}. 빛-물질 상호작용의 맥락에서 \(\Keff\)는 광자와 원자 집단 사이의 간섭성을 측정한다.
	
	여기서 우리는 음의 시간 지연 실험이 YUCT의 직접적인 결과임을 보여준다. 그러나 우리는 \(\tau_g\)와 \(\Keff\) 사이의 선형 관계를 가정했던 이전의 시도를 수정해야 한다. YUCT의 정보이론적 구조로부터 유도된 올바른 관계는 로그적이다. 이는 예측된 온도 의존성을 극적으로 변화시켜 현재의 실험적 정밀도로는 무시할 수 있게 만든다. 따라서 우리는 이전에 발표된 13\% 예측을 철회하고 수정된 자기-일관된 분석을 제공한다.
	
	\section{YUCT의 핵심 개념}
	
	\subsection{조정 효율 \(\Keff\)}
	
	YUCT에서 임의의 상호작용 시스템은 사전 \(D\) (공유 프로토콜, 상태)와 특정 조정된 행동을 활성화하는 인덱스 \(i\)에 의해 특성화된다. 조정의 효율은
	\begin{equation}
		\Keff = \frac{H(\text{행동})}{H(\text{인덱스})},
	\end{equation}
	여기서 \(H\)는 섀넌 엔트로피를 나타낸다. 양자 시스템의 경우 \(\Keff\)는 매우 커질 수 있으며, 완벽한 얽힘의 극한에서 무한대에 접근한다.
	
	\subsection{범용 스케일링 법칙}
	
	변동과 오차는 식~\eqref{eq:scaling}을 따른다. 이 법칙은 사전의 프랙탈 기하학으로부터 발생하며, 중성자 붕괴에서 은하 회전 곡선에 이르기까지 다양한 시스템에서 경험적으로 확인되었다 \cite{Yakushev2026c}. 로그 형태는 필수적이다: 이는 인덱스 공간에서의 유효 분해능이 사전 크기에 따라 로그적으로만 증가한다는 사실을 반영한다.
	
	\subsection{\(\Keff\)의 온도 의존성}
	
	많은 물리적 시스템에서 조정 효율은 열적 변동으로 인해 온도가 증가함에 따라 감소한다. 임계 현상과 프랙탈 스케일링으로부터 우리는
	\begin{equation}
		\Keff(T) = K_0 \left(\frac{T}{T_0}\right)^{-\gamma}, \quad \gamma > 0,
		\label{eq:K_T}
	\end{equation}
	을 가지며, 여기서 \(\gamma\)는 재료 의존적 지수이다. 곱 \(\betac\gamma\)는 지구물리학적 데이터로부터 독립적으로 결정되었다: \(\betagamma = 0.20 \pm 0.03\) \cite{Yakushev2026b}. 이는 \(\gamma = \betagamma / \betac = 0.20 / (2/3) = 0.30\)을 의미한다. 이 값은 무질서 시스템에 대한 일반적인 스케일링 논증과 일치한다.
	
	\section{광자-원자 상호작용의 YUCT 모델}
	
	냉각 원자 집단에 입사하는 공명 광자 펄스를 고려하자. 원자는 \textbf{사전} \(D\)로 작용한다: 그들은 특정 양자 상태로 준비된다 (예: 레이저 냉각 및 트래핑을 통해). 광자는 \textbf{인덱스} – 그의 파동 묶음 모양과 주파수 – 를 운반한다. 상호작용은 광자가 투과, 흡수 또는 산란될 수 있는 조정된 상태를 초래한다.
	
	Angulo 등 \cite{Angulo2024}의 실험에서 관심 대상은 투과된 광자에 대한 \textbf{원자 여기 시간} \(\tau_T\)이다. YUCT로부터, 이 시간은 군 지연 \(\tau_g\)에 비례한다. Thompson 등 \cite{Thompson2023}의 약값 분석은 관련 영역에서 투과된 광자에 대해 \(\tau_T = \tau_g\)임을 보여준다.
	
	결정적 단계는 \(\tau_g\)를 \(\Keff\)와 연결하는 것이다. 군 지연은 광자와 원자 집단 사이의 간섭성 상호작용으로부터 발생한다. 이 상호작용의 크기는 \(\Keff\) 자체 (엄청나게 클 수 있음)에 비례하는 것이 아니라, 광자가 원자 상태에 대해 운반하는 \textbf{정보}의 양에 비례한다. YUCT에서 관련 정보 측도는 조정 효율의 로그인데, 이는 인덱스 길이 \(\ell = \log_2 M\)가 구별 가능한 상태의 수를 결정하기 때문이다. 따라서 우리는 다음과 같이 가정한다:
	\begin{equation}
		|\tau_g| = \alpha_0 \, \ln\Keff,
		\label{eq:tau_log}
	\end{equation}
	여기서 \(\alpha_0\)는 시간의 차원을 가진 기본 시간 상수이다. 루비듐 D2 선에 대해, 자연 선폭 \(\Gamma = 2\pi \times 6\) MHz는 특성 시간 \(\tau_{\text{nat}} = 1/\Gamma \approx 26\) ns을 제공한다. \(K_0 \sim 10^4\)에서 실험값 \(|\tau_g| \approx 30\) ns에 피팅하면 \(\alpha_0 = 30\,\text{ns} / \ln(10^4) = 30 / 9.21 \approx 3.26\) ns을 얻는다. 이는 타당하다: 기본 상호작용 시간은 수 나노초 정도이다.
	
	\subsection{군 지연의 온도 의존성}
	
	식~\eqref{eq:K_T}을 식~\eqref{eq:tau_log}에 대입하면
	\begin{equation}
		|\tau_g(T)| = \alpha_0 \ln\!\left(K_0 \left(\frac{T}{T_0}\right)^{-\gamma}\right)
		= \alpha_0 \ln K_0 - \alpha_0 \gamma \ln\!\left(\frac{T}{T_0}\right).
	\end{equation}
	\(|\tau_g(T_0)| = \alpha_0 \ln K_0\)로 표기하면, 우리는
	\begin{equation}
		|\tau_g(T)| = |\tau_g(T_0)| - \alpha_0 \gamma \ln\!\left(\frac{T}{T_0}\right).
		\label{eq:tau_T_log}
	\end{equation}
	을 얻는다. 이것이 핵심 결과이다: 군 지연의 크기는 온도 증가에 따라 \textbf{로그적으로} 감소한다. 온도 비 \(T/T_0 = 2\)에 대해 변화는
	\begin{equation}
		\Delta |\tau_g| = \alpha_0 \gamma \ln 2.
	\end{equation}
	\(\alpha_0 \approx 3.26\) ns와 \(\gamma = 0.30\)을 사용하면, \(\Delta |\tau_g| \approx 3.26 \times 0.30 \times 0.693 \approx 0.68\) ns을 찾는다. 기준값 \(|\tau_g(T_0)| \approx 30\) ns에 대한 상대적 변화는 \(0.68 / 30 \approx 2.3\%\)이다. 이는 여전히 실험의 통계적 불확실성 (\(\pm 0.22\) 단위 \(\tau_0\), 즉 약 \(\pm 5.7\) ns) 내에 있다. 그러나 이 효과는 검출 가능성의 경계에 있다. 더 중요하게, 변화는 \textbf{로그적}이며, 거듭제곱 법칙이 아니고, 그 크기는 이전에 주장된 13\%보다 훨씬 작다.
	
	\subsection{이전의 (잘못된) 선형 모델과의 비교}
	
	이전 연구에서는 선형 관계 \(|\tau_g| \propto \Keff\)가 가정되었다. 이는 \(|\tau_g| \propto T^{-\gamma}\)로 이어졌고, \(\gamma = 0.30\)으로 온도 두 배 시 13\%의 변화를 초래했다. 이 예측은 양자 시스템에서 \(\Keff\)의 역할에 대한 오해에 기반했다. 올바른 로그 관계는 \(\Keff\)가 엔트로피의 비율이지 직접적인 물리적 장 진폭이 아니기 때문에 발생한다. 선형 모델은 YUCT의 정보이론적 기초와 양립할 수 없으며, 이에 따라 철회된다.
	
	\section{기존 데이터로부터의 수치적 일관성 검사}
	
	Angulo 등 \cite{Angulo2024}의 실험은 \(\alpha_0\)와 \(K_0\)를 추정할 수 있게 하는 특정 데이터 점을 제공한다. 10 ns 펄스, 광학 깊이 OD~4 구성에 대해 보고된 원자 여기 시간은 \(\tau_T = -1.14 \tau_0\)였다. \(\tau_0 \approx 26\) ns를 취하면, \(|\tau_g| \approx 30\) ns을 얻는다. \(\Keff(T_0) = K_0\)가 \(10^4\) 정도라고 가정하면 (전형적인 냉각 원자 집단은 \(10^6\)개의 원자와 높은 간섭성을 가짐), \(\alpha_0 = 30\,\text{ns} / \ln(10^4) \approx 3.26\) ns을 얻는다. 이 값은 기본 광자-원자 상호작용의 특성 시간 척도로서 합리적이다. 추가적인 자유 매개변수는 필요하지 않다.
	
	\section{표준 양자 광학과의 비교}
	
	표준 양자 광학 형식론에서 공명 펄스의 군 지연 \(\tau_g\)는 투과 위상의 주파수에 대한 미분, \(\tau_g = d\phi/d\omega\)로 주어진다. 로렌츠 원자 응답에 대해, 이는 이상 분산으로 인해 공명의 날개에서 음수가 될 수 있는 특성 형태를 제공한다 \cite{Milonni2004}. 그러나 표준 이론은 도플러 넓힘 (온도에서 선형, 로그 아님)과 같은 사소한 효과를 제외하고, \(\tau_g\)의 원자 구름 온도에 대한 의존성을 예측하지 않는다. 관찰된 음의 시간 지연은 에너지 교환을 포함하지 않는 간섭성 간섭 효과로 이해된다.
	
	YUCT는 지연의 크기를 조정 효율 \(\Keff\)와 연결함으로써 표준 이론을 넘어선다. 로그 의존성은 특징적인 신호이다: 원자 밀도, 광학 깊이, 또는 펄스 지속 시간을 변경하여 \(\Keff\)를 변화시킬 수 있다면, 지연은 \(\ln\Keff\)로 스케일되어야 한다. 이는 온도 변화에 의존하지 않는 검증 가능한 예측이다.
	
	\section{수정된 실험적 제안}
	
	온도 의존성이 극도로 약하기 때문에 (로그적, 온도 두 배 시 상대적 변화가 몇 퍼센트에 불과), 현재 정밀도로는 검출될 가능성이 낮다. 따라서 우리는 \(\Keff\)를 더 통제된 방식으로 직접 변화시키는 대체 실험을 제안한다.
	
	\subsection{광학 깊이의 변화}
	
	조정 효율 \(\Keff\)는 집단 내 원자 수와 상호작용 강도에 비례한다. 광학 깊이를 줄임으로써 (예: 원자 밀도를 낮춤), \(\Keff\)가 감소한다. 군 지연 크기는 \(|\tau_g| = \alpha_0 \ln\Keff\)을 따라야 한다. \(\Keff\)를 10배 줄이면 (예: \(10^4\)에서 \(10^3\)으로), \(|\tau_g|\)의 변화는 \(\alpha_0 (\ln 10^4 - \ln 10^3) = \alpha_0 \ln 10 \approx 3.26 \times 2.3026 \approx 7.5\) ns가 되며, 이는 기준값의 약 25\%이고 쉽게 측정 가능하다. 이러한 실험은 간단하다: 다른 원자 밀도에 대해 \(\tau_T\) 측정을 반복하고 (다른 모든 매개변수는 고정) 로그 스케일링을 검증한다.
	
	\subsection{펄스 지속 시간의 변화}
	
	광자 펄스의 정보 내용 (인덱스 길이)은 시간적 지속 시간과 대역폭에 의존한다. 더 짧은 펄스는 더 많은 스펙트럼 정보를 운반하므로 더 큰 \(\Keff\)를 가진다. 펄스 지속 시간을 1 ns에서 20 ns까지 체계적으로 변화시킴으로써, 군 지연이 \(\ln(\text{대역폭})\)에 따라 어떻게 스케일하는지 관찰할 수 있다. 이는 로그 법칙의 또 다른 직접적 검증을 제공한다.
	
	\subsection{온도 의존성을 영점 검증으로}
	
	예측된 온도 효과는 작지만, 정확히 0은 아니다. 영점 측정 (실험 오차 내에서 검출 가능한 변화가 없음)은 우리의 로그 모델과 일치할 것이며, 이전의 선형 모델은 명확한 13\% 변화를 예측했을 것이다. 따라서 기존 데이터는 이미 선형 모델을 지지하지 않으며, 미래의 고정밀 실험은 궁극적으로 로그 이동을 검출할 수 있을 것이다. 우리는 실험자들이 둘 다 검증할 것을 권장한다.
	
	\section{광자 질량 예측과의 연결: 일관성 검사}
	\label{sec:photon_mass_consistency}
	
	YUCT 프레임워크는 또한 동일한 질량 사다리 \(m = m_e q^{N_f}\) (\(q=(3/2)^{1/3}\))로부터 유도된 광자의 정지 질량에 대한 맹목적 예측을 산출했다. 실험적 상한과 일치하는 최적 노드는 \(N_\gamma = -410\)이며, 이는
	\[
	m_\gamma \approx 7.3 \times 10^{-19}~\text{eV} \quad (\text{부록 Photon Mass 참조})
	\]
	을 제공한다.
	
	이러한 미세한 비-영 광자 질량이 냉각 원자에서의 군 지연 측정에 영향을 미칠 수 있는지 의문이 들 수 있다. \(m_\gamma\)에 해당하는 콤프턴 파장은
	\[
	\lambda_c = \frac{\hbar}{m_\gamma c} \approx 2.7\times10^{11}~\text{m},
	\]
	이며, 이는 원자 구름 (전형적 크기 \(\sim 1\) mm)보다 훨씬 크다. 따라서 유한 광자 질량으로 인한 전파 효과는 이 실험에서 완전히 무시할 수 있다. 음의 시간 지연은 순수한 조정 효과이며, 절대 질량 척도와 무관하다.
	
	반대로, 단일 프레임워크가 다음을 모두 올바르게 예측한다는 사실은:
	\begin{itemize}
		\item \(\ln\Keff\)에 의해 크기가 설정된 음의 군 지연의 존재 (본 부록),
		\item 광자 질량의 상한 (\(m_\gamma < 10^{-18}\) eV) 및 이산 노드 \(N_\gamma = -410\) (부록 Photon Mass),
	\end{itemize}
	강력한 일관성 검사를 구성한다. 두 예측은 매우 다른 물리적 영역 (양자 광학 대 기본 입자 특성)에 관련되지만, 동일한 기본 상수 \(q\), \(S_{\mathrm{odd}}/S_{\mathrm{even}}\), 그리고 조정 효율 \(\Keff\)의 개념을 공유한다. 이들을 호환 가능하게 만들기 위해 특별한 조정이 필요하지 않다.
	
	따라서 음의 광자 지연 실험은 YUCT 광자 질량 예측과 모순되지 않으며; 대신 둘 다 통합 조정 패러다임을 지지하는 독립적 기둥으로 작용한다.
	
	\section{이전의 맹목적 예측 철회}
	
	이 부록의 이전 버전 (V.2.0)에서 우리는 군 지연이 \(T^{-0.20}\)으로 스케일할 것이라고 잘못 예측했으며, 이는 온도 두 배 시 13\% 변화로 이어졌다. 이 예측은 잘못된 선형 관계 \(|\tau_g| \propto \Keff\)에 기반했다. 여기서 유도된 올바른 로그 관계는 온도 효과가 약 두 자릿수 더 작음을 보여준다. 따라서 13\% 효과에 대한 주장은 무효이며 공식적으로 철회된다. 이 간과에 대해 사과드리며, 불일치를 식별해준 수학적 패치의 공동 저자들에게 감사드린다.
	
	그럼에도 불구하고, YUCT의 핵심 개념 – 음의 군 지연이 조정 효율의 발현이라는 것 – 은 여전히 유효하다. 음의 원자 여기 시간의 실험적 확인은 이미 광자와 원자가 조정된 사전을 공유한다는 아이디어를 지지한다. 수정된 스케일링 법칙은 이제 미래 검증을 위한 정제된 이론적 틀을 제공한다.
	
	\section{결론}
	
	우리는 음의 광자 지연 실험에 대한 YUCT 분석을 수정했다. YUCT의 범용 오차 법칙은 로그적이며, 이에 따라 군 지연과 조정 효율 사이의 관계는 로그적이어야 한다: \(|\tau_g| = \alpha_0 \ln\Keff\). 이는 매우 약한 온도 의존성 (로그적, 온도 두 배 시 상대적 변화가 몇 퍼센트 정도)으로 이어지며, 이전에 주장된 13\% 효과를 제거한다. 잘못된 예측은 철회된다. 광학 깊이나 펄스 지속 시간의 변화와 같은 대체 검증은 로그 스케일링에 대해 훨씬 더 민감한 탐침을 제공한다. 음의 시간 지연의 존재는 양자 상호작용의 조정된 본질에 대한 놀라운 확인으로 남아 있으며, YUCT는 일관된 정보이론적 해석을 제공한다.
	
	\section*{감사의 말}
	
	저자는 선형 가정과 YUCT의 로그 구조 사이의 불일치를 지적해준 수학적 패치의 공동 저자들에게 감사드린다. 또한 YUCT 세미나 참가자들에게 자극적인 토론에, 그리고 토론토 그룹 (Aephraim Steinberg, Daniela Angulo 및 동료들)에게 그들의 아름다운 실험에 감사드린다.
	
	\section*{데이터 가용성}
	
	모든 계산 및 추가 자료는 \url{https://github.com/Alexey-Yakushev-YUCT/YPSDC}에서 이용 가능하다.
	
	\begin{thebibliography}{99}
		
		\bibitem{Angulo2024}
		D. Angulo, K. Thompson, A. Jiao, H. Wiseman, A. Steinberg, ``Experimental evidence that a photon can spend a negative amount of time in an atom cloud,'' arXiv:2409.03680 [quant-ph] (2024).
		
		\bibitem{Nixon2024}
		V.-M. Nixon et al., ``Measuring the time transmitted photons spend as atomic excitations,'' American Physical Society DAMOP Conference, Session Y07 (June 2024).
		
		\bibitem{Yakushev2026a}
		A. V. Yakushev, ``Yakushev Unified Coordination Theory (YUCT),'' Zenodo, \url{https://doi.org/10.5281/zenodo.18444599} (2026).
		
		\bibitem{Yakushev2026b}
		A. V. Yakushev, ``Appendix P: Temperature Dependence of Gravity,'' Zenodo (2026).
		
		\bibitem{Yakushev2026c}
		A. V. Yakushev, ``Appendix L: Fractal Error Scaling,'' Zenodo (2026).
		
		\bibitem{Yakushev2026d}
		A. V. Yakushev, ``Appendix X: Generalization of Shannon Information Theory,'' Zenodo (2026).
		
		\bibitem{Thompson2023}
		K. Thompson et al., ``How much time does a photon spend as an atomic excitation before being transmitted?'' arXiv:2310.00432 (2023).
		
		\bibitem{Brillouin1960}
		L. Brillouin, \textit{Wave Propagation and Group Velocity} (Academic Press, 1960).
		
		\bibitem{Milonni2004}
		P. Milonni, \textit{Fast Light, Slow Light and Left-Handed Light} (CRC Press, 2004).
		
		\bibitem{Wiseman2023}
		H. M. Wiseman, A. M. Steinberg, and M. Hallaji, ``Obtaining a single-photon weak value from experiments using a strong coherent state,'' AVS Quantum Sci. \textbf{5}, 024401 (2023).
		
	\end{thebibliography}
	
\end{document}